\section{G2: Shock fitting and the post-shock flow}
\label{sec:g2}
\subsection{Problem and governing equations}
In a TOP bell the curvature discontinuity at N generates compressions that coalesce
into an internal shock upstream of the last left-running characteristic, so the shock
affects the wall \cite[PDF p.~26, printed p.~18, Sec.~3.5]{ostlund}. The same shock
underlies the cap-shock plume pattern \cite[PDF p.~3, printed p.~480]{frey}. The
repository detects the fold and stops. No open numeric TOP internal-shock location or
strength data were found: Östlund gives the mechanism only, and Ruf et al.\ show the
shock only in CFD images \cite{ruf2010}.

For a stationary discontinuity with normal and tangential velocity components
$u_n$ and $u_t$,
\begin{equation}
\rho_1u_{n1}=\rho_2u_{n2},\quad p_1+\rho_1u_{n1}^2=p_2+\rho_2u_{n2}^2,\quad
h_1+\tfrac12u_{n1}^2=h_2+\tfrac12u_{n2}^2,\quad u_{t1}=u_{t2}
\label{eq:rh}
\end{equation}
\cite[PDF p.~104, Eqs.~3.38--3.40; PDF pp.~150--151, Eqs.~4.3--4.6]{anderson}.

\textbf{Closures and branches.}
\begin{itemize}
\item \textbf{Calorically perfect:} the $\theta$--$\beta$--$M$ relation with weak and
strong branches \cite[PDF pp.~152--154, printed pp.~136--138]{anderson}. Its entropy
relations (PDF p.~109, Eqs.~3.61--3.63) do not carry over to a changing composition.
\item \textbf{General normal shock:} iterate on $\rho_1/\rho_2$
\cite[PDF pp.~665--666, printed pp.~649--650, Eqs.~17.1--17.10]{anderson}, or apply
Newton's method to $\ln T$ and $\ln p$
\cite[PDF pp.~39--40, printed pp.~33--34, Eqs.~7.6--7.19]{gordon1994}.
\item \textbf{Equilibrium oblique shock:} Moeckel prescribes the shock angle
\cite[PDF pp.~3--4, printed pp.~2--3, Eqs.~5--7]{moeckel1957}. No source selects the
branch, so S0 applies Eq.~\eqref{eq:rh} to the normal component and takes the
smallest $\beta>\mu_1$ giving the required deflection. This is our derivation,
tested against the calorically perfect limit.
\end{itemize}

\textbf{Post-shock flow.} Behind a curved shock the flow is rotational, so the MOC
must carry entropy and total enthalpy along streamlines
\cite[PDF pp.~423--424, Sec.~11.9]{anderson}. Shyne gives the rotational
characteristic equations with a streamline family
\cite[PDF p.~33, printed p.~28, Eqs.~58--62]{shyne1988}\cite[PDF pp.~8--9]{shyne1990}.
Fitting and capturing are distinguished in \cite[PDF p.~438, Sec.~11.15]{anderson}.

\textbf{Shock onset.} In one-dimensional unsteady piston compression, same-family
characteristics form an envelope, and a shock must occur on or before it
\cite[Sec.~16, p.~40; Sec.~24, pp.~58--59]{courant1944}. Inouye, Rakich and Lomax
start a shock where two same-family Mach lines intersect, at their mean slope
\cite[PDF pp.~12--14, printed pp.~10--12]{inouye1965}.

\textbf{Evidence (E10, E10b; Fig.~\ref{fig:e10}).} On the current-chart bell the fold
occurs on the first $C^-$ from the bell after N, at a grazing angle that shrinks with
refinement. $r/R_t$ stays within 0.941--0.950 while $x$ drifts from 55.3 to 50.7\,mm:
far more stable than $x$, but not shown to converge. Straight-characteristic envelope
estimates do not locate the fold, because it forms through curved $C^-$. The
upper-bound strength estimate is $p_{02}/p_{01}=0.917$. E10 remains INCOMPLETE.

\subsection{Approaches and choice}
\begingroup\footnotesize
\par\addvspace{\LTpre}{\centering\begin{tabular}{@{}P{0.28\linewidth}P{0.53\linewidth}P{0.12\linewidth}@{}}
\toprule\textbf{Approach} & \textbf{Assessment} & \textbf{Verdict}\\\midrule
Detect, refuse, map crossings & Locates the end of the smooth solution; the first crossing bounds the onset & S0a\\
Weak-shock isentropic continuation & Entropy error third order in strength; only a policy for very weak coalescences (MOC-31) & policy\\
Fitted shock with rotational downstream & Keeps MOC accuracy and cost; verifiable against exact jumps & \textbf{S1}\\
In-repo capturing FV & General; large self-verification burden & reject\\
External capturing CFD & Independent check of the location & CFD1\\
\bottomrule
\end{tabular}\par}\addvspace{\LTpost}
\endgroup
S0a locates the envelope from the net's own crossing locus. S0 is a standalone jump
operator on the T0 equation of state. S1 fits one right-running internal shock with a
supersonic downstream region. Reflections at the axis, Mach disks and cap shocks are
refused (S2). Figure~\ref{fig:shock_topology} shows the proposed topology.

% Generated by builds/moc_extensions_plan_2026-11-21/writer_t/gen_shock_topology.py (schematic geometry).
\begin{figure}[htbp]
\centering
\begin{tikzpicture}[x=1in,y=1in,font=\sffamily\scriptsize,
  chm/.style={draw=ink,line width=0.45pt},
  chp/.style={draw=teal,line width=0.8pt},
  sl/.style={draw=muted,line width=0.45pt,dash pattern=on 2.2pt off 1.3pt},
  shk/.style={draw=accent,line width=1.5pt},
  env/.style={draw=accent,line width=0.7pt,dash pattern=on 0.5pt off 1.2pt},
  lead/.style={draw=muted,line width=0.3pt},
  lbl/.style={text=ink,inner sep=1pt,align=left},
  pO/.style={circle,draw=ink,fill=white,inner sep=0pt,minimum size=4.4pt,line width=0.6pt},
  pF/.style={circle,draw=ink,fill=ink,inner sep=0pt,minimum size=4.4pt},
  pS/.style={rectangle,draw=ink,fill=white,inner sep=0pt,minimum size=4.2pt,line width=0.6pt},
  pD/.style={diamond,draw=accent,fill=white,inner sep=0pt,minimum size=6pt,line width=0.7pt},
  pX/.style={draw=accent,line width=1.1pt}]
% kernel: start line, throat arc to N, C- from N to the axis
\fill[pale] (0.740,0.000) .. controls (0.710,0.240) and (0.530,0.430) .. (0.550,0.553) arc[start angle=-84.26,end angle=-56.00,radius=0.5in] .. controls (1.330,0.335) and (1.700,0.100) .. (2.250,0.000) -- cycle;
% region 2: behind the fitted front, below the dividing streamline through C
\fill[accent!8] (1.550,0.700) -- (1.591,0.692) -- (1.631,0.684) -- (1.672,0.676) -- (1.712,0.668) -- (1.753,0.660) -- (1.794,0.651) -- (1.834,0.643) -- (1.875,0.634) -- (1.915,0.625) -- (1.956,0.616) -- (1.997,0.607) -- (2.037,0.598) -- (2.078,0.589) -- (2.118,0.579) -- (2.159,0.569) -- (2.199,0.560) -- (2.240,0.550) -- (2.281,0.540) -- (2.321,0.530) -- (2.362,0.519) -- (2.402,0.509) -- (2.443,0.498) -- (2.484,0.488) -- (2.524,0.477) -- (2.565,0.466) -- (2.605,0.455) -- (2.646,0.444) -- (2.687,0.432) -- (2.727,0.421) -- (2.768,0.409) -- (2.808,0.398) -- (2.849,0.386) -- (2.890,0.374) -- (2.930,0.362) -- (2.971,0.349) -- (3.011,0.337) -- (3.052,0.324) -- (3.092,0.312) -- (3.133,0.299) -- (3.174,0.286) -- (3.214,0.273) -- (3.255,0.260) -- (3.295,0.247) -- (3.336,0.233) -- (3.377,0.219) -- (3.417,0.206) -- (3.458,0.192) -- (3.498,0.178) -- (3.539,0.164) -- (3.580,0.150) -- (3.620,0.135) -- (3.661,0.121) -- (3.701,0.106) -- (3.742,0.091) -- (3.783,0.076) -- (3.823,0.061) -- (3.864,0.046) -- (3.904,0.031) -- (3.945,0.016) -- (3.985,0.000) -- (4.900,0) -- (4.900,1.311) -- (4.844,1.306) -- (4.788,1.300) -- (4.732,1.294) -- (4.677,1.287) -- (4.621,1.281) -- (4.565,1.275) -- (4.509,1.268) -- (4.453,1.262) -- (4.398,1.255) -- (4.342,1.248) -- (4.286,1.241) -- (4.230,1.234) -- (4.174,1.227) -- (4.118,1.220) -- (4.063,1.213) -- (4.007,1.205) -- (3.951,1.197) -- (3.895,1.190) -- (3.839,1.182) -- (3.783,1.174) -- (3.728,1.166) -- (3.672,1.157) -- (3.616,1.149) -- (3.560,1.140) -- (3.504,1.132) -- (3.448,1.123) -- (3.393,1.114) -- (3.337,1.104) -- (3.281,1.095) -- (3.225,1.086) -- (3.169,1.076) -- (3.113,1.066) -- (3.058,1.056) -- (3.002,1.046) -- (2.946,1.035) -- (2.890,1.025) -- (2.834,1.014) -- (2.778,1.003) -- (2.723,0.992) -- (2.667,0.981) -- (2.611,0.969) -- (2.555,0.957) -- (2.499,0.945) -- (2.443,0.933) -- (2.388,0.921) -- (2.332,0.908) -- (2.276,0.895) -- (2.220,0.882) -- (2.164,0.868) -- (2.108,0.854) -- (2.053,0.840) -- (1.997,0.826) -- (1.941,0.812) -- (1.885,0.797) -- (1.829,0.781) -- (1.773,0.766) -- (1.718,0.750) -- (1.662,0.734) -- (1.606,0.717) -- (1.550,0.700) -- cycle;
% axis, wall, exit plane
\draw[muted,line width=0.4pt,dash pattern=on 6pt off 1.5pt on 1pt off 1.5pt] (0.02,0) -- (4.980,0);
\draw[ink,line width=0.9pt] (0.182,0.682) arc[start angle=-135,end angle=-90,radius=0.45in];
\draw[ink,line width=0.9pt] (0.500,0.550) arc[start angle=-90,end angle=-56.00,radius=0.5in];
\draw[ink,line width=0.9pt] (0.780,0.635) .. controls (1.699,1.256) and (3.073,1.711) .. (4.900,2.000);
\draw[muted,line width=0.4pt,dash pattern=on 2pt off 1.5pt] (4.900,0) -- (4.900,2.000);
\node[text=muted,anchor=south,inner sep=1pt] at (4.900,2.000) {exit};
\node[text=muted,anchor=north west,inner sep=1pt] at (0.02,-0.01) {axis};
% transonic start line
\draw[teal,line width=1.0pt] (0.740,0.000) .. controls (0.710,0.240) and (0.530,0.430) .. (0.550,0.553);
% kernel boundary: C- from N to the axis
\draw[chm] (0.780,0.635) .. controls (1.330,0.335) and (1.700,0.100) .. (2.250,0.000);
% C- from the throat arc to the axis point Ka
\draw[chm] (0.621,0.565) .. controls (0.921,0.325) and (1.200,0.07) .. (1.550,0.000);
% C+ leaving the axis point: interior point, crosses the front, ends at a wall point in region 2
\draw[chp] (1.550,0.000) -- (2.511,0.480) -- (4.140,1.865);
% wall-emitted compression C-, tangent to the envelope; they end on the fitted front
\draw[chm] (1.056,0.809) -- (1.550,0.700);
\draw[chm] (1.095,0.831) -- (1.617,0.687);
\draw[chm] (1.142,0.858) -- (1.715,0.668);
\draw[chm] (1.196,0.888) -- (1.819,0.646);
\draw[chm] (1.271,0.929) -- (1.950,0.618);
% envelope near the cusp: locus of first same-family crossings (S0a); the front lies on or before it
\draw[env] (1.550,0.700) -- (1.585,0.692) -- (1.621,0.684) -- (1.656,0.675) -- (1.692,0.666) -- (1.727,0.657) -- (1.762,0.647) -- (1.798,0.637) -- (1.833,0.626) -- (1.869,0.616) -- (1.904,0.605) -- (1.940,0.593) -- (1.975,0.581) -- (2.010,0.569) -- (2.046,0.556) -- (2.081,0.544) -- (2.117,0.530) -- (2.152,0.517) -- (2.188,0.503) -- (2.223,0.489) -- (2.258,0.474) -- (2.294,0.459) -- (2.329,0.444) -- (2.365,0.428) -- (2.400,0.412);
% fitted front (S1); dashed where it would reach the axis
\draw[shk] (1.550,0.700) -- (1.602,0.690) -- (1.655,0.680) -- (1.707,0.669) -- (1.759,0.658) -- (1.812,0.648) -- (1.864,0.636) -- (1.917,0.625) -- (1.969,0.613) -- (2.021,0.602) -- (2.074,0.590) -- (2.126,0.577) -- (2.178,0.565) -- (2.231,0.552) -- (2.283,0.539) -- (2.335,0.526) -- (2.388,0.513) -- (2.440,0.499) -- (2.493,0.485) -- (2.545,0.471) -- (2.597,0.457) -- (2.650,0.443) -- (2.702,0.428) -- (2.754,0.413) -- (2.807,0.398) -- (2.859,0.383) -- (2.911,0.367) -- (2.964,0.351) -- (3.016,0.335) -- (3.069,0.319) -- (3.121,0.303) -- (3.173,0.286) -- (3.226,0.269) -- (3.278,0.252) -- (3.330,0.235) -- (3.383,0.217) -- (3.435,0.200) -- (3.487,0.182) -- (3.540,0.164) -- (3.592,0.145) -- (3.645,0.127);
\draw[accent,line width=0.9pt,dash pattern=on 3pt off 2pt] (3.645,0.127) -- (3.679,0.114) -- (3.713,0.102) -- (3.747,0.090) -- (3.781,0.077) -- (3.815,0.064) -- (3.849,0.052) -- (3.883,0.039) -- (3.917,0.026) -- (3.951,0.013) -- (3.985,0.000);
% streamlines: one crosses the front and kinks toward the axis; dividing streamline through C
\draw[sl] (0.800,0.175) -- (0.834,0.181) -- (0.868,0.187) -- (0.903,0.193) -- (0.937,0.199) -- (0.971,0.205) -- (1.005,0.210) -- (1.039,0.216) -- (1.073,0.221) -- (1.108,0.226) -- (1.142,0.232) -- (1.176,0.237) -- (1.210,0.242) -- (1.244,0.247) -- (1.278,0.252) -- (1.312,0.257) -- (1.347,0.261) -- (1.381,0.266) -- (1.415,0.271) -- (1.449,0.275) -- (1.483,0.280) -- (1.518,0.284) -- (1.552,0.288) -- (1.586,0.293) -- (1.620,0.297) -- (1.654,0.301) -- (1.688,0.305) -- (1.723,0.309) -- (1.757,0.313) -- (1.791,0.317) -- (1.825,0.321) -- (1.859,0.325) -- (1.893,0.329) -- (1.928,0.333) -- (1.962,0.336) -- (1.996,0.340) -- (2.030,0.344) -- (2.064,0.347) -- (2.098,0.351) -- (2.133,0.354) -- (2.167,0.358) -- (2.201,0.361) -- (2.235,0.365) -- (2.269,0.368) -- (2.303,0.371) -- (2.338,0.374) -- (2.372,0.378) -- (2.406,0.381) -- (2.440,0.384) -- (2.474,0.387) -- (2.508,0.390) -- (2.543,0.393) -- (2.577,0.396) -- (2.611,0.399) -- (2.645,0.402) -- (2.679,0.405) -- (2.713,0.408) -- (2.748,0.411) -- (2.782,0.413) -- (2.816,0.411) -- (2.850,0.408) -- (2.884,0.406) -- (2.918,0.403) -- (2.953,0.400) -- (2.987,0.397) -- (3.021,0.394) -- (3.055,0.394) -- (3.089,0.397) -- (3.123,0.399) -- (3.158,0.401) -- (3.192,0.403) -- (3.226,0.406) -- (3.260,0.408) -- (3.294,0.410) -- (3.328,0.412) -- (3.363,0.414) -- (3.397,0.416) -- (3.431,0.418) -- (3.465,0.421) -- (3.499,0.423) -- (3.533,0.425) -- (3.567,0.427) -- (3.602,0.428) -- (3.636,0.430) -- (3.670,0.432) -- (3.704,0.434) -- (3.738,0.436) -- (3.773,0.438) -- (3.807,0.440) -- (3.841,0.442) -- (3.875,0.443) -- (3.909,0.445) -- (3.943,0.447) -- (3.978,0.449) -- (4.012,0.451) -- (4.046,0.452) -- (4.080,0.454) -- (4.114,0.456) -- (4.148,0.457) -- (4.183,0.459) -- (4.217,0.461) -- (4.251,0.462) -- (4.285,0.464) -- (4.319,0.465) -- (4.353,0.467) -- (4.388,0.469) -- (4.422,0.470) -- (4.456,0.472) -- (4.490,0.473) -- (4.524,0.475) -- (4.558,0.476) -- (4.593,0.478) -- (4.627,0.479) -- (4.661,0.480) -- (4.695,0.482) -- (4.729,0.483) -- (4.763,0.485) -- (4.798,0.486) -- (4.832,0.487) -- (4.866,0.489) -- (4.900,0.490);
\draw[sl] (1.00,0.509) -- (1.550,0.700) -- (1.606,0.717) -- (1.662,0.734) -- (1.718,0.750) -- (1.773,0.766) -- (1.829,0.781) -- (1.885,0.797) -- (1.941,0.812) -- (1.997,0.826) -- (2.053,0.840) -- (2.108,0.854) -- (2.164,0.868) -- (2.220,0.882) -- (2.276,0.895) -- (2.332,0.908) -- (2.388,0.921) -- (2.443,0.933) -- (2.499,0.945) -- (2.555,0.957) -- (2.611,0.969) -- (2.667,0.981) -- (2.723,0.992) -- (2.778,1.003) -- (2.834,1.014) -- (2.890,1.025) -- (2.946,1.035) -- (3.002,1.046) -- (3.058,1.056) -- (3.113,1.066) -- (3.169,1.076) -- (3.225,1.086) -- (3.281,1.095) -- (3.337,1.104) -- (3.393,1.114) -- (3.448,1.123) -- (3.504,1.132) -- (3.560,1.140) -- (3.616,1.149) -- (3.672,1.157) -- (3.728,1.166) -- (3.783,1.174) -- (3.839,1.182) -- (3.895,1.190) -- (3.951,1.197) -- (4.007,1.205) -- (4.063,1.213) -- (4.118,1.220) -- (4.174,1.227) -- (4.230,1.234) -- (4.286,1.241) -- (4.342,1.248) -- (4.398,1.255) -- (4.453,1.262) -- (4.509,1.268) -- (4.565,1.275) -- (4.621,1.281) -- (4.677,1.287) -- (4.732,1.294) -- (4.788,1.300) -- (4.844,1.306) -- (4.900,1.311);
% unit processes
\draw[-{Stealth[length=1.3mm]},accent,line width=0.5pt] (3.415,0.254) -- (3.486,0.462) node[anchor=south west,inner sep=0.5pt,text=accent] {$\mathbf{n}$};
\node[pO] at (3.387,0.169) {};
\node[pF,draw=accent,fill=accent] at (3.415,0.254) {};
\node[pO] at (1.797,0.123) {};
\node[pS] at (1.550,0.000) {};
\node[pF] at (4.140,1.865) {};
\node[pD] at (2.511,0.480) {};
\node[circle,fill=accent,inner sep=0pt,minimum size=4.6pt] at (1.550,0.700) {};
\fill[ink] (0.780,0.635) circle (1.4pt);
\draw[pX] (3.920,-0.065) -- (4.050,0.065) (3.920,0.065) -- (4.050,-0.065);
% labels
\node[lbl,anchor=south east,inner sep=0.5pt] at (0.760,0.665) {N};
\draw[lead] (0.720,0.755) -- (0.55,1.18) node[lbl,anchor=south,align=center] {curvature change at N:\\throat arc $\to$ TOP/TIC bell};
\node[lbl,text=teal,align=center] at (0.27,0.27) {transonic\\start line};
\node[lbl,text=muted] at (1.02,0.09) {kernel};
\node[lbl,anchor=north east,inner sep=0.5pt,text=accent] at (1.520,0.680) {C};
\draw[lead] (1.296,0.878) -- (1.75,1.62) node[lbl,anchor=south,align=center] {wall-emitted compressions ($C^-$) coalesce;\\envelope cusp C, where the front starts};
\node[lbl,anchor=west] at (2.42,0.10) {\textbf{1} isentropic};
\node[lbl,anchor=west,align=left] at (3.55,0.86) {\textbf{2} rotational, $s=s(\psi)$:\\varies across streamlines};
\node[lbl,anchor=south west,text=accent,font=\sffamily\scriptsize\bfseries] at (4.045,0.03) {S2};
% legend: line styles
\draw[chm] (5.120,2.050) -- (5.380,2.050); \node[lbl,anchor=north west,inner sep=0pt] at (5.420,2.105) {$C^-$ (right-running)};
\draw[chp] (5.120,1.895) -- (5.380,1.895); \node[lbl,anchor=north west,inner sep=0pt] at (5.420,1.950) {$C^+$ (left-running)};
\draw[sl] (5.120,1.740) -- (5.380,1.740); \node[lbl,anchor=north west,inner sep=0pt] at (5.420,1.795) {streamline: third\\characteristic in 2 \citep{shyne1988}};
\draw[env] (5.120,1.430) -- (5.380,1.430); \node[lbl,anchor=north west,inner sep=0pt] at (5.420,1.485) {envelope locus (S0a)};
\draw[shk] (5.120,1.275) -- (5.380,1.275); \node[lbl,anchor=north west,inner sep=0pt] at (5.420,1.330) {fitted front (S1)};
% legend: unit processes
\node[pO] at (5.250,1.070) {};
\node[lbl,anchor=north west,inner sep=0pt] at (5.420,1.125) {interior point};
\node[pF] at (5.250,0.915) {};
\node[lbl,anchor=north west,inner sep=0pt] at (5.420,0.970) {wall point};
\node[pS] at (5.250,0.760) {};
\node[lbl,anchor=north west,inner sep=0pt] at (5.420,0.815) {axis point};
\node[pO] at (5.200,0.605) {}; \node[pF,draw=accent,fill=accent] at (5.300,0.605) {};
\node[lbl,anchor=north west,inner sep=0pt] at (5.420,0.660) {shock point: state pair\\1\,|\,2, S0 jump along $\mathbf{n}$};
\node[pD] at (5.250,0.295) {};
\node[lbl,anchor=north west,inner sep=0pt] at (5.420,0.350) {shock--characteristic\\intersection};
\draw[pX] (5.200,-0.065) -- (5.300,0.035) (5.200,0.035) -- (5.300,-0.065);
\node[lbl,anchor=north west,inner sep=0pt] at (5.420,0.040) {refused: axis reflection,\\Mach disk (S2)};
\end{tikzpicture}
\caption{Proposed shock-fitted topology in a TOP bell; a schematic, not a computed solution. Compressions emitted by the wall after the curvature change at N coalesce. S0a locates their envelope from the net's own first same-family crossings, on or before which a shock must form \citep{courant1944}; S1 starts one right-running front at the cusp C and fits it with the S0 general-EOS jump. Behind the front the flow is rotational and the streamline is a third characteristic \citep{shyne1988}; a front that would reach the axis or the wall is refused (S2 scope).}
\label{fig:shock_topology}
\end{figure}


\begingroup\footnotesize
\par\addvspace{\LTpre}{\centering\begin{tabular}{@{}P{0.09\linewidth}P{0.36\linewidth}P{0.30\linewidth}P{0.17\linewidth}@{}}
\toprule\textbf{ID} & \textbf{Numerical verification} & \textbf{Benchmark agreement} & \textbf{Physical validation}\\\midrule
S0a & Reproduces t01 (done); crossing-locus refinement; weak-coalescence classes declared before testing & CFD1 captured-shock location & none\\
S0 & $\theta$--$\beta$--$M$ to $10^{-10}$; normal-shock limit on the same EOS & Moeckel equilibrium air (after T2) & none\\
S1 & Weak-shock limit; conservation across the front; grid convergence of the location & CFD1 location within the captured thickness & none\\
S2 & Refusal tests & n/a & n/a\\
\bottomrule
\end{tabular}\par}\addvspace{\LTpost}
\endgroup
